\documentclass{article}
\usepackage{graphicx}
\usepackage{amssymb}
\usepackage{amsmath}
\usepackage{stmaryrd}
\title{Solution to the uniqueness of Dunford decomposition}
\author{Luigi de Alfaro}
\date{2024}

\begin{document}

\maketitle

\section{Problem Statement}

Let \( n \in \mathbb{N} \), and let \( N, N', S, S' \in M_n(\mathbb{R}) \) such that:
\begin{itemize}
    \item \( S \) and \( N \) commute (\( SN = NS \)),
    \item \( S' \) and \( N' \) commute (\( S'N' = N'S' \)),
    \item \( S \) and \( S' \) are diagonalizable,
    \item \( N \) and \( N' \) are nilpotent,
    \item \( S + N = S' + N' \).
\end{itemize}

Prove that \( S = S' \) and \( N = N' \).

Rather than the usual proof, we will solve this problem considering \( P(S + N) \) for a well-chosen polynomial \( P \in \mathbb{R}[X] \). The polynomial is obtained by Hermite interpolation: we prescribe not only its values at the eigenvalues of \( S \), but also its first \( n-1 \) derivatives there.

I stumbled across this problem as a nineteen-year-old and decided to give it a go, finding this solution by myself. At the time, I had no knowledge of the Dunford decomposition. I am aware that this problem has been solved for a long time, but I wanted to share the proof I found.

\section{Notation}


Since \( S \) is diagonalizable in \( M_n(\mathbb{R}) \), its eigenvalues are real. We write
\( \lambda_1, \dots, \lambda_r \) for its \emph{distinct} eigenvalues, and
\( d_1, \dots, d_n \) for its eigenvalues repeated according to multiplicity, so that the
\( d_j \) are exactly the diagonal entries of a diagonalization of \( S \) and
\( \{d_1, \dots, d_n\} = \{\lambda_1, \dots, \lambda_r\} \) as sets. Note also that
\( N^n = 0 \), since a nilpotent matrix of \( M_n(\mathbb{R}) \) has nilpotency index at
most \( n \).

\section{Solution}

\subsection{Step 1: Spectral properties of \( S \) and \( S' \)}

The matrices \( S \) and \( N \) commute, and both have a characteristic polynomial that
splits over \( \mathbb{R} \): that of \( S \) because \( S \) is diagonalizable in
\( M_n(\mathbb{R}) \), and that of \( N \) because \( N \) is nilpotent, so
\( \chi_N = X^n \). Two commuting matrices with split characteristic polynomials are
simultaneously trigonalizable, so there is a basis in which \( S \) and \( N \) are both
upper triangular.

In such a basis, the diagonal of \( S \) carries its eigenvalues \( d_1, \dots, d_n \),
repeated according to multiplicity. Since \( N \) is nilpotent, in such a basis, \( N \) is strictly upper triangular. Consequently \( S + N \) is upper
triangular with the same diagonal as \( S \), and therefore:

\( S \) and \( S + N \) have the same eigenvalues, with the same
multiplicities. The same argument gives that \( S' \) and \( S' + N' \) have the same eigenvalues, with the same
multiplicities. Since
\( S + N = S' + N' \), we conclude that  \( S \) and \( S' \)
have the same eigenvalues, with the same multiplicities.

\subsection{Step 2: Polynomial representation}

Let \( P \in \mathbb{R}[X] \) be a polynomial:
\[
P(X) = \sum_{k=0}^d a_k X^k.
\]
Since \( S \) is diagonalizable, it can be written as:
\[
S = Q 
\begin{pmatrix}
d_1 & & \\
& \ddots & \\
& & d_n
\end{pmatrix}
Q^{-1},
\]
where \( d_1, \dots, d_n \) are the eigenvalues of \( S \) repeated according to
multiplicity.
Using the commutativity of \( S \) and \( N \), we apply the binomial theorem to \( P(S + N) \):
\[
P(S + N) = \sum_{k=0}^d a_k \sum_{i=0}^k \binom{k}{i} Q 
\begin{pmatrix}
d_1^{k-i} & & \\
& \ddots & \\
& & d_n^{k-i}
\end{pmatrix}
Q^{-1} N^i.
\]
Rearranging the sums gives:
\[
P(S + N) = \sum_{i=0}^d \frac{1}{i!} Q 
\begin{pmatrix}
\sum_{k=i}^d a_k \frac{k!}{(k-i)!} d_1^{k-i} & & \\
& \ddots & \\
& & \sum_{k=i}^d a_k \frac{k!}{(k-i)!} d_n^{k-i}
\end{pmatrix}
Q^{-1} N^i,
\]
which simplifies to:
\[
P(S + N) = \sum_{i=0}^d \frac{1}{i!} Q 
\begin{pmatrix}
P^{(i)}(d_1) & & \\
& \ddots & \\
& & P^{(i)}(d_n)
\end{pmatrix}
Q^{-1} N^i.
\]

\subsection{Step 3: Choosing the polynomial \( P \)}

Suppose \( P \) is chosen such that:
\begin{itemize}
    \item \( P(\lambda_j) = \lambda_j \) for all \( j \in \llbracket 1, r \rrbracket \),
    \item \( P^{(i)}(\lambda_j) = 0 \) for all \( j \in \llbracket 1, r \rrbracket \) and all \( i \in \llbracket 1, n-1 \rrbracket \).
\end{itemize}
% CORRIGE : les conditions sont imposees aux valeurs propres distinctes ; comme chaque
% d_m est l'un des lambda_j, elles valent automatiquement pour tous les d_m.
Since every \( d_m \) is one of the \( \lambda_j \), these conditions give
\( P(d_m) = d_m \) and \( P^{(i)}(d_m) = 0 \) for \( 1 \le i \le n-1 \) and every
\( m \in \llbracket 1, n \rrbracket \).

In the expansion of Step 2, the term \( i = 0 \) is then
\( Q\,\mathrm{diag}(d_1, \dots, d_n)\,Q^{-1} = S \); the terms with
\( 1 \le i \le n-1 \) vanish because the diagonal matrix is zero; and the terms with
\( i \ge n \) vanish because \( N^i = 0 \). Hence
\[
  P(S+N) = S .
\]
The same polynomial works for \( S' \), since \( S \) and \( S' \) have the same
eigenvalues by Step 1, so \( P(S' + N') = S' \). As \( S + N = S' + N' \), it follows
that \( S = S' \).

\subsection{Step 4: Existence of the polynomial \( P \)}

The conditions of Step 3 prescribe, at each of the \( r \) distinct points
\( \lambda_1, \dots, \lambda_r \), the value of \( P \) and its first \( n-1 \)
derivatives. This is exactly a Hermite interpolation problem, and it always has a
solution.

Consider the linear map
\[
  \Phi : \mathbb{R}_{rn-1}[X] \longrightarrow \mathbb{R}^{rn}, \qquad
  P \longmapsto \Big( P(\lambda_j), P'(\lambda_j), \dots, P^{(n-1)}(\lambda_j) \Big)_{1 \le j \le r},
\]
where \( \mathbb{R}_{rn-1}[X] \) denotes the polynomials of degree strictly less than
\( rn \). Both spaces have dimension \( rn \), so it suffices to show that \( \Phi \) is
injective.

Let \( P \in \ker \Phi \). For each \( j \), the first \( n \) derivatives of \( P \) at
\( \lambda_j \) (orders \( 0 \) to \( n-1 \)) vanish, so \( \lambda_j \) is a root of
\( P \) of multiplicity at least \( n \). The \( \lambda_j \) being pairwise distinct, the
polynomial \( \prod_{j=1}^{r} (X - \lambda_j)^n \), of degree \( rn \), divides \( P \).
Since \( \deg P < rn \), this forces \( P = 0 \).

Hence \( \Phi \) is bijective, and there exists a polynomial \( P \) of degree less than
\( rn \) satisfying all the conditions of Step 3.

\subsection{Step 5: Conclusion}

Applying \( P(S + N) = S \) and \( P(S' + N') = S' \), and given \( S + N = S' + N' \), we
obtain \( S = S' \). Subtracting from \( S + N = S' + N' \) then gives \( N = N' \).

\end{document}
